% Generated from locale/tr/content/turing-machines/machines-computations/church-turing-thesis.tex; do not hand-edit.


\olfileid[tr]{tur}{mac}{ctt}
\olsection{Church--Turing Tezi}

Turing makinelerinin, etkili yöntem kavramının kesin bir karşılığı olması
amaçlanır. Turing, hem etkili yöntem kavramını hem de Turing makinesi kavramını
kavrayan herkesin, etkili bir yöntemle yapılabilen her şeyin bir Turing
makinesiyle de yapılabileceği sezgisine sahip olacağını düşünüyordu. Etkili
yöntem kavramı için önerilen bütün öteki kesin karşılıkların Turing makinesi
kavramına kaplamsal olarak eşdeğer çıkması, yani tam olarak aynı fonksiyonlar
kümesini hesaplayabilmeleri bu iddiayı destekler. Bu iddiaya
\emph{Church--Turing tezi} denir.

\begin{defn}[Church--Turing tezi]
\emph{Church--Turing Tezi}, etkili bir yöntemle hesaplanabilen her şeyin
Turing hesaplanabilir olduğunu söyler.
\end{defn}

Church--Turing tezine iki biçimde başvurulur. İlk kullanım biçimi tembellik
için bir mazerettir. Bir şeyi hesaplayan etkili bir yöntemin, diyelim ki
``sözde kod'' ile verilmiş bir betimine sahip olduğumuzu varsayalım. O zaman
onu gerçekten kurmamış olsak bile aynı fonksiyonun bir Turing makinesi
tarafından hesaplandığı iddiasını gerekçelendirmek için Church--Turing tezine
başvurabiliriz.

Church--Turing tezinin öteki kullanımı felsefi bakımdan daha ilginçtir.
Turing makinelerince hesaplanamayan fonksiyonların bulunduğu gösterilebilir.
Buradan Church--Turing tezi kullanılarak bu fonksiyonların hiçbir yöntemle
etkili biçimde hesaplanamayacağı sonucuna varılabilir. Çünkü böyle bir yöntem
olsaydı Church--Turing tezine göre onun için bir Turing makinesi de bulunurdu.
Dolayısıyla bir fonksiyonu hesaplayan hiçbir Turing makinesinin bulunmadığını
kanıtlayabilirsek etkili bir yöntem de bulunamaz. Özellikle Church--Turing
tezine, durma problemi denen problemin yalnızca Turing makinelerince değil,
hiçbir biçimde etkili olarak çözülemeyeceğini ileri sürmek için başvurulur.

